% Bagian: first-order-logic
% Bab: axiomatic-deduction
% Subbagian: provability

% Verifikasi sifat-sifat keterbuktian yang diperlukan untuk himpunan
% konsisten maksimal dalam bab kelengkapan.

\documentclass[../../../include/open-logic-section]{subfiles}

\begin{document}

\iftag{FOL}
      {\olfileid[id]{fol}{axd}{prv}}
      {\olfileid[id]{pl}{axd}{prv}}

\olsection{Sifat-Sifat \usetoken{S}{derivability}}

\begin{prop}[Monotonisitas] Jika $\Gamma \subseteq \Delta$ dan
$\Gamma \Proves !A$, maka $\Delta \Proves !A$ juga.
\end{prop}

\begin{proof} Setiap $\Gamma_0 \subseteq \Gamma$ yang berhingga juga
merupakan subhimpunan berhingga dari~$\Delta$, sehingga !!a{derivation} atas
$\Gamma_0 \Proves !A$ juga menunjukkan $\Delta \Proves !A$.
\end{proof}

\begin{prop}\ollabel{prop:provability} $\Gamma \Proves !A$ jika dan hanya
jika $\Gamma\cup \{\lnot!A \}$ tak konsisten.
\end{prop}

\begin{prop} \begin{enumerate}
\item \ollabel{prop:provability-contr} Jika $\Gamma \Proves !A$ dan
$\Gamma \cup \{ !A\} \Proves \lfalse$, maka $\Gamma$ tak konsisten.

\item \ollabel{prop:provability-lnot} Jika $\Gamma \cup \{!A\}
\Proves \lfalse$, maka $\Gamma \Proves \lnot !A$.

\item \ollabel{prop:provability-exhaustive} Jika $\Gamma \cup \{!A\}
\Proves \lfalse$ dan $\Gamma \cup \{\lnot !A\} \Proves
\lfalse$, maka $\Gamma \Proves \lfalse$.

\tagitem{prvOr}{\ollabel{prop:provability-lor-left} Jika
$\Gamma \cup \{!A\} \Proves \lfalse$ dan $\Gamma \cup \{!B\}
\Proves \lfalse$, maka $\Gamma \cup \{!A \lor !B\} \Proves \lfalse$.}{}

\tagitem{prvOr}{\ollabel{prop:provability-lor-right} Jika
$\Gamma \Proves !A$ atau $\Gamma \Proves !B$, maka
$\Gamma \Proves !A \lor !B$.}{}

\tagitem{prvAnd}{\ollabel{prop:provability-land-left} Jika
$\Gamma \Proves !A \land !B$, maka $\Gamma \Proves !A$ dan
$\Gamma \Proves !B$.}{}

\tagitem{prvAnd}{\ollabel{prop:provability-land-right} Jika
$\Gamma \Proves !A$ dan $\Gamma \Proves !B$, maka
$\Gamma \Proves !A \land !B$.}{}

\tagitem{prvIf}{\ollabel{prop:provability-mp} Jika $\Gamma \Proves !A$
dan $\Gamma \Proves !A \lif !B$, maka $\Gamma \Proves !B$.}{}

\tagitem{prvIf}{\ollabel{prop:provability-lif} Jika
$\Gamma \Proves \lnot !A$ atau $\Gamma \Proves !B$, maka
$\Gamma \Proves !A \lif !B$.}{}
\end{enumerate} \end{prop}

\begin{proof}
\begin{enumerate}
\item Andaikan $\Gamma_0 \cup \{!A\} \Proves \lfalse$.

\begin{derivation}
1. & $\Gamma_0 \cup \{!A\} \Proves \lfalse$ & hip. \\
2. & $\Gamma_1 \Proves !A$ & hip.\\
3. & $\Gamma_0 \cup \Gamma_1 \Proves \lfalse$ & Cut \\
\end{derivation}

Karena $\Gamma_0 \subseteq \Gamma$ dan $\Gamma_1 \subseteq \Gamma$,
$\Gamma_0 \cup \Gamma_1 \subseteq \Gamma$; jadi,
$\Gamma \Proves \lfalse$.

\item Andaikan $\Gamma_0 \cup\{!A\} \Proves \lfalse$.

\begin{derivation}
1. & $\Gamma_0 \cup\{!A\} \Proves \lfalse$ & hip. \\
2. & $\Gamma_0 \Proves !A \lif \lfalse$ & Teorema Deduksi\\
3. & $\Gamma_0 \Proves (!A \lif \lfalse) \lif \lnot !A$ & Ax0\\
4. & $\Gamma_0 \Proves \lnot !A$ & MP 2, 3 \\
\end{derivation}

\item Andaikan $\Gamma_0 \cup \{ !A \} \Proves \lfalse$ dan
$\Gamma_1 \cup \{\lnot !A \} \Proves \lfalse$.

\begin{derivation}
1. & $\Gamma_0 \cup\{!A\} \Proves \lfalse$ & hip. \\
2. & $\Gamma_0 \Proves !A \lif \lfalse$ & Teorema Deduksi \\
3. & $\Gamma_1 \cup \{\lnot !A \} \Proves \lfalse$ & hip.\\
4. & $\Gamma_0 \Proves (!A \lif \lfalse) \lif \lnot !A$ & Ax0\\
5. & $\Gamma_0 \Proves \lnot !A$ & MP 2, 4\\
6. & $\Gamma_0 \cup \Gamma_1 \Proves \lfalse$ & Cut 3, 5 \\
\end{derivation}

Karena $\Gamma_0 \subseteq \Gamma$ dan $\Gamma_1 \subseteq \Gamma$,
$\Gamma_0 \cup \Gamma_1 \subseteq \Gamma$. Jadi,
$\Gamma \Proves \lfalse$.

% prop:provability-lor-left
\tagitem{defOr}{}{
\iftag{probOr}{Latihan.}{Latihan.}}

% prop:provability-lor-right
\tagitem{defOr}{}{ \iftag{probOr}{Latihan.}{
Andaikan $\Gamma_0 \Proves !A$.

\begin{derivation}
1. & $\Gamma_0 \Proves !A$ & hip. \\
2. & $\Gamma_0 \Proves !A \lif (!A \lor !B)$ & Ax0\\
3. & $\Gamma_0 \Proves !A \lor !B$ & MP 1, 2 \\
\end{derivation}

Jadi, $\Gamma \Proves !A \lor !B$. Pembuktian untuk kasus ketika
$\Gamma \Proves !B$ serupa.}}

% prop:provability-land-left
\tagitem{defAnd}{}{ \iftag{probAnd}{Latihan}{
Andaikan $\Gamma_0 \Proves !A \land !B$.

\begin{derivation}
1. & $\Gamma_0 \Proves !A \land !B$ & hip. \\
2. & $\Gamma_0 \Proves (!A\land !B) \lif !A $ & Ax0\\
3. & $\Gamma_0 \Proves !A$ & MP 1, 2 \\
\end{derivation}

Jadi, $\Gamma \Proves !A$. !!a{derivation} yang serupa menunjukkan bahwa
$\Gamma \Proves !B$.}}

% prop:provability-land-right
\tagitem{defAnd}{}{\iftag{probAnd}{Latihan.}{Latihan.}}

% prop:provability-mp
\tagitem{defIf}{}{ \iftag{probIf}{Latihan.}{
Andaikan $\Gamma_0 \Proves !A$ dan $\Gamma_1 \Proves !A \lif !B$.

\begin{derivation}
1. & $\Gamma_0 \Proves !A$ & hip. \\
2. & $\Gamma_1 \Proves !A \lif !B $ & hip.\\
3. & $\Gamma_0 \cup \Gamma_1 \Proves !B$ & MP 1, 2\\
\end{derivation}

Karena $\Gamma_0 \cup \Gamma_1 \subseteq \Gamma$, ini berarti bahwa
$\Gamma \Proves !B$.}}

% prop:provability-lif
\tagitem{defIf}{}{\iftag{probIf}{Latihan.}{Pertama-tama andaikan
$\Gamma \Proves \lnot !A$.

\begin{derivation}
1. & $\Gamma_0 \Proves \lnot !A$ & hip. \\
2. & $\Gamma_0 \Proves \lnot !A \lif (!A \lif !B) $ & Ax0\\
3. & $\Gamma_0 \Proves !A \lif !B$ & MP 1, 2\\
\end{derivation}

Pembuktian untuk $\Gamma \Proves !B$ serupa:

\begin{derivation}
1. & $\Gamma_0 \Proves !B$ & hip. \\
2. & $\Gamma_0 \Proves !B \lif (!A \lif !B) $ & Ax0\\
3. & $\Gamma_0 \Proves !A \lif !B$ & MP 1, 2\\
\end{derivation}}}

\end{enumerate}
\end{proof}

\begin{probtag}{probOr,probAnd,probIf} Lengkapi pembuktian
\olref[fol][axd][prv]{prop:provability}.
\end{probtag}

\begin{thm}[Generalisasi] \ollabel{thm:Generalization} Jika
$\Gamma \Proves !A$ dan $x$ tidak bebas dalam !!{formula} mana pun dalam
$\Gamma$, maka $\Gamma \Proves \lforall[x][!A]$.
\end{thm}

\begin{proof} Dengan induksi pada $\mathsf{Thm}_\Gamma$: jika $!A$ merupakan
aksioma, maka $\lforall[x][!A]$ juga merupakan aksioma. Jika
$!A\in \Gamma$, maka $x$ tidak bebas dalam $!A$, sehingga
$!A \lif \lforall[x][!A]$ merupakan aksioma (\textbf{Ax3}), dan melalui MP
kita memperoleh $\Gamma \Proves \lforall[x][!A]$. Sekarang andaikan $!A$
diperoleh melalui MP karena $\Gamma \Proves !B$ dan
$\Gamma \Proves !B \lif !A$. Menurut hipotesis induksi,
$\Gamma \Proves \lforall[x][!B]$ dan
$\Gamma \Proves \lforall[x][( !B \lif !A)]$. Namun, menurut
\textbf{Ax2} kita juga mempunyai
\[
\Gamma \Proves
\lforall[x][( !B \lif !A)] \lif (\lforall[x][ !B] \lif \lforall[x][!A]),
\]
dan dua penerapan MP memberikan $\Gamma \Proves \lforall[x][!A]$, seperti
yang diinginkan.
\end{proof}

\begin{thm}\ollabel{thm:WeakGen} (\emph{Generalisasi Lemah atas Konstanta})
Jika $\Gamma \Proves !A$ dan $c$ merupakan !!a{constant} yang tidak muncul
dalam $\Gamma$, maka terdapat variabel~$x$ yang tidak muncul dalam $!A$
sedemikian sehingga $\Gamma \Proves \lforall[x][\Subst{!A}{x}{c}]$, dan
pembuktiannya tidak melibatkan~$c$.
\end{thm}

\begin{proof} Misalkan $!A_1,\ldots,!A_n$ merupakan pembuktian $!A$ dari
$\Gamma$, sehingga $!A=!A_n$. Pilih variabel~$x$ yang tidak muncul dalam
$!A_1,\ldots,!A_n$, dan tinjau barisan baru
$\Subst{!A_1}{x}{c},\ldots,\Subst{!A_n}{x}{c}$. Barisan semacam itu merupakan
pembuktian $\Subst{!A}{x}{c}$ dari $\Gamma$. Bahkan, bagi setiap
$i=1,\ldots,n$:
\begin{itemize}
\item jika $!A_i$ merupakan aksioma, maka $\Subst{!A_i}{x}{c}$ juga
merupakan aksioma;
\item jika $!A_i \in \Gamma$, maka, karena $c$ tidak muncul dalam
$\Gamma$, kita mempunyai $\Subst{!A_i}{x}{c} = !A_i \in \Gamma$;
\item jika $!A_i$ diperoleh dari $!A_j$ dan $!A_j \lif !A_i$, maka
$\Subst{!A_i}{x}{c}$ diperoleh melalui MP dari $\Subst{!A_j}{x}{c}$ dan
$\Subst{(!A_j \lif !A_i)}{x}{c}
= \Subst{!A_j}{x}{c} \lif \Subst{!A_i}{x}{c}$.
\end{itemize}
Jelas bahwa !!{constant}~$c$ tidak lagi muncul dalam barisan baru tersebut.
Sekarang misalkan $\Gamma'$ terdiri atas !!{formula} dalam $\Gamma$ yang
muncul dalam !!{derivation} atas $\Subst{!A}{x}{c}$; maka $x$ tidak bebas
dalam !!{formula} mana pun dalam $\Gamma'$ dan
$\Gamma' \Proves \Subst{!A}{x}{c}$. Menurut generalisasi
(\olref{thm:Generalization}),
$\Gamma' \Proves \lforall[x][\Subst{!A}{x}{c}]$; jadi, menurut
monotonisitas, $\Gamma \Proves \lforall[x][\Subst{!A}{x}{c}]$, seperti
yang diinginkan.
\end{proof}

\begin{explain}
Tujuan kita ialah mengganti syarat dalam \olref{thm:WeakGen} bahwa $x$ sama
sekali tidak muncul dalam $!A$ dengan syarat yang lebih lemah bahwa $x$ tidak
bebas dalam $!A$ dan !!{free for} $c$ dalam $!A$. Jelas bahwa hal ini dapat
dicapai dengan mengganti !!{variable} terikat---jadi, itulah yang mula-mula
akan kita buktikan.
\end{explain}

\begin{lem}
\ollabel{lem:free-for}
Jika $x$ bebas untuk $c$ dalam $!A$ dan $y$ tidak bebas dalam $!A$, maka
$x$ !!{free for} $y$ dalam $\Subst{!A}{y}{c}$.
\end{lem}

\begin{proof} Jika $x$ tidak !!{free for} $y$ dalam
$\Subst{!A}{y}{c}$, maka suatu kemunculan bebas $y$ dalam
$\Subst{!A}{y}{c}$ berada dalam cakupan kuantor~$\lforall[x]$. Namun,
menurut hipotesis, $y$ tidak bebas dalam $!A$, sehingga semua kemunculan
semacam itu berasal dari substitusi~$\Subst{}{y}{c}$. Jadi, suatu kemunculan
$c$ berada dalam cakupan $\lforall[x]$, dan $x$ tidak !!{free for} $c$
dalam $!A$.
\end{proof}

\begin{lem}[Penggantian Variabel Terikat]
\ollabel{lem:ChangeBdVar}
Jika $x$ dan $y$ tidak bebas dalam $!A$ dan keduanya !!{free for} $c$
dalam $!A$, maka $\Proves \lforall[x][\Subst{!A}{x}{c}] \equiv
\lforall[y][\Subst{!A}{y}{c}]$ (dan pembuktiannya tidak melibatkan $c$).
\end{lem}

\begin{proof}
Dari hipotesis, $y$ !!{free for} $c$ dalam $!A$ dan $x$ tidak bebas dalam
$!A$; jadi, menurut \olref{lem:free-for} sebelumnya, $y$ !!{free for} $x$
dalam $\Subst{!A}{x}{c}$. Akibatnya,
$\lforall[x][\Subst{!A}{x}{c} \lif \Subst{!A}{x}{c} \Subst{}{y}{x}]$
merupakan aksioma (\textbf{Ax1}). Karena $x$ tidak bebas dalam $!A$, kita
mempunyai
$\Subst{!A}{x}{c} \Subst{}{y}{x} = \Subst{!A}{y}{c}$, sehingga
\[
\Proves \lforall[x][\Subst{!A}{x}{c}] \lif \Subst{!A}{y}{c},
\]
dan menurut Teorema Deduksi,
$\lforall[x][\Subst{!A}{x}{c}] \Proves \Subst{!A}{y}{c}$. Karena $y$
tidak bebas dalam $!A$, variabel tersebut juga tidak bebas dalam
$\lforall[x][\Subst{!A}{x}{c}]$; jadi, menurut Generalisasi,
$\lforall[x][\Subst{!A}{x}{c}] \Proves
\lforall[y][\Subst{!A}{y}{c}]$, dan Teorema Deduksi kembali memberikan
$\Proves \lforall[x][\Subst{!A}{x}{c}] \lif
\lforall[y][\Subst{!A}{y}{c}]$. Pembuktian
$\Proves \lforall[y][\Subst{!A}{y}{c}] \lif
\lforall[x][\Subst{!A}{x}{c}]$ sepenuhnya simetris, sehingga kesimpulan
mengikuti menurut Taut.
\end{proof}

\begin{thm}[Generalisasi Kuat atas Konstanta]
\ollabel{thm:StrongGen}
Jika $\Gamma \Proves !A$, !!{constant}~$c$ tidak muncul dalam $\Gamma$,
dan $x$ tidak bebas dalam $!A$ tetapi !!{free for} $c$ dalam $!A$, maka
$\Gamma \Proves \lforall[x][\Subst{!A}{x}{c}]$.
\end{thm}

\begin{proof}
Karena $\Gamma \Proves !A$ dan $c$ tidak muncul dalam $\Gamma$, menurut
Generalisasi Lemah terdapat !!a{variable}~$y$ yang tidak muncul dalam $!A$
sedemikian sehingga $\Gamma \Proves \lforall[y][\Subst{!A}{y}{c}]$.
Karena $y$ sama sekali tidak muncul dalam $!A$, variabel tersebut tidak bebas dalam $!A$
dan juga !!{free for} $c$ dalam $!A$. Jika, selain itu (menurut hipotesis),
$x$ tidak bebas dalam $!A$ dan !!{free for} $c$ dalam $!A$, maka syarat
untuk penggantian !!{variable} terikat terpenuhi, sehingga
$\Proves \lforall[x][\Subst{!A}{x}{c}] \equiv
\lforall[y][\Subst{!A}{y}{c}]$; akibatnya,
$\Gamma \Proves \lforall[x][\Subst{!A}{x}{c}]$.
\end{proof}

\begin{thm}
\ollabel{thm:provability-quantifiers}
\begin{tagenumerate}{prvEx,prvAll}
\tagitem{prvEx}{Jika $\Gamma \Proves !A(t)$, maka $\Gamma \Proves
  \lexists[x][!A(x)]$.}{}
\tagitem{prvAll}{Jika $\Gamma \Proves \lforall[x][!A(x)]$, maka $\Gamma
  \Proves !A(t)$.}{}
\end{tagenumerate}
\end{thm}

\begin{proof}
\begin{tagenumerate}{prvEx,prvAll}
%% Memerlukan aksioma bagi kuantor eksistensial
\tagitem{prvEx}{Latihan.}{}
\tagitem{prvAll}{Andaikan $\Gamma_0 \Proves \lforall[x][!A(x)]$.

\begin{derivation}
1. & $\Gamma_0 \Proves \lforall[x][!A(x)]$ & hip. \\
2. & $\Gamma_0 \Proves \lforall[x][!A(x)] \lif \Subst{!A}{t}{x}$ & Ax1 \\
3. & $\Gamma_0 \Proves \Subst{!A}{t}{x}$ & MP 1, 2 \\
\end{derivation}}{}

\end{tagenumerate}
\end{proof}

\end{document}
